%! Choosing a Sample: Four Sampling Techniques — Unit 1, Lesson 1.3 — Lesson Plan
% PILOT SHAPE (2026-09-06): modeled on AP Statistics lesson 1.4 minus its AP Practice page.
% GRADUAL-RELEASE plan (warm-up / guided notes and practice / spoken debrief / close and START
% THE HOMEWORK). The guided notes carry the release: I do (two sections) -> we do (Guided
% Practice). THERE IS NO INDEPENDENT PRACTICE SET IN THE NOTES — the homework is the individual
% practice, and the period ends with students starting it. Teacher-facing. Breakable boxes are
% `skillbox` with a \boxguard ahead; the two tabularx boxes are `fixedskillbox` (never splits).
% No group activity, no tier boxes, no exit ticket.
\documentclass[10pt]{article}
\usepackage{saar-article}
\usepackage{saar-boxes}
\usepackage{graphicx}   % -article does NOT load graphicx
\graphicspath{{images/}}
\newcommand{\TallMath}[1]{$\displaystyle #1\rule[-1.4em]{0pt}{3.2em}$}
\newcommand{\UnitNumberName}{Unit 1: Asking Questions with Data: Study Design \quad}
\newcommand{\LessonNumberName}{Lesson 1.3: Choosing a Sample --- Four Sampling Techniques}

\begin{document}
\begin{center}
    {\Large\bfseries \CourseName \\
    {\small\bfseries \UnitNumberName \LessonNumberName}}
\end{center}

\begin{tcolorbox}[colback=frost, colframe=cerulean, boxrule=0.9pt, arc=2mm,
    left=2mm, right=2mm, top=1.5mm, bottom=1.5mm]
    \textbf{Primary Objective:} Students will name the sampling frame and say whether chance
    or convenience chose a sample; carry out a simple random sample from generator digits and
    a systematic sample from a random start; compute a stratified allocation in proportion to
    the strata and count what a cluster sample reaches; tell a stratum from a cluster; and
    choose the technique a context calls for with a reason that names \emph{how the groups are
    built}.
    \\[2pt]
    \textbf{Standards:} PS.DC.2a, PS.DC.2b; AFDA.DA.2b \quad$\cdot$\quad PS.DC.1d carried
    forward in the population/sample items \quad$\cdot$\quad previews PS.DC.2c and
    AFDA.DA.2e--g (Lesson 1.4)
    \\[2pt]
    \textbf{Lesson model:} \emph{Gradual release --- I do, we do, you do --- all of it inside
    the guided notes.} There is no group activity. The notes name the vocabulary as they teach
    it and deliver the instruction in two sections; the Guided Practice is worked together with
    students holding the pen; \textbf{the homework is the individual practice}, and students
    start it in class while you circulate. The whole-class debrief is \emph{spoken} --- there is
    no practice set at the end of the notes and no exit ticket.
    \\[2pt]
    \textbf{Note on the homeroom map:} the crux depends on Riverbend's homerooms being grouped
    by grade ($1$--$9$ ninth, $10$--$17$ tenth, $18$--$24$ eleventh, $25$--$30$ twelfth). Do
    not reveal that until the second half of section~2 --- the hook and the first half of
    section~2 run on the hat draw alone.
\end{tcolorbox}

\renewcommand{\arraystretch}{1.4}
\boxguard
\begin{skillbox}[Learning Targets \& Key Understandings]{goldbox}
\begin{tabularx}{\linewidth}{@{}X|X@{}}
\textbf{Learning targets (student language)} & \textbf{Key understandings (the ``why'')} \\[2pt]
\hline
\vspace{-6pt}
\begin{itemize}[leftmargin=1.1em, nosep, after=\vspace{2pt}]
  \item I can name the \textbf{sampling frame} and say whether \textbf{chance} or
        \textbf{convenience} chose a sample, by asking who could never have been picked.
        \hfill \textit{PS.DC.2a}
  \item I can carry out a \textbf{simple random sample} from generator digits and a
        \textbf{systematic sample} with $k = N \div n$. \hfill \textit{PS.DC.2a; AFDA.DA.2b}
  \item I can compute a \textbf{stratified} allocation in proportion and count what a
        \textbf{cluster sample} reaches. \hfill \textit{PS.DC.2a; AFDA.DA.2b}
  \item I can tell a \textbf{stratum} from a \textbf{cluster} and choose --- with a reason ---
        the technique a context calls for. \hfill \textit{PS.DC.2b; AFDA.DA.2b}
\end{itemize}
&
\vspace{-6pt}
\begin{itemize}[leftmargin=1.1em, nosep, after=\vspace{2pt}]
  \item \emph{The method is decided before any data exist}, and it decides what the statistic
        is worth. Lesson 1.2 said a statistic estimates a parameter; whether it estimates it
        \emph{well} starts here.
  \item \textbf{Target misconception:} that ``chance chose'' settles it --- that a random draw
        of groups is as good as a random draw of individuals, so stratified and cluster are
        interchangeable ways of ``using the groups.'' They are not: a cluster sample works only
        when each cluster is a \emph{small mix of the whole}. Maya's two random homerooms are
        $60$ ninth graders and no seniors, because a homeroom is alike inside --- a stratum.
        \emph{Alike inside: some from every group. A mix of the whole: everyone from a few.}
  \item \textbf{Second misconception:} that \emph{random} means \emph{unplanned}. The teacher
        who surveys her own third period picked no favorites, yet a fourth-period student could
        never have been chosen --- convenience, however fair it felt.
  \item Nobody left off the \textbf{frame} can ever be sampled, however good the method.
  \item ``Best'' is not a property of a technique; it is a property of a technique \emph{in a
        context} --- what list exists, and how the groups are built (PS.DC.2b).
\end{itemize}
\end{tabularx}
\end{skillbox}

\boxguard[24]
\begin{skillbox}[{Vocabulary, Concepts, \& Theorems --- taught directly in the guided notes}]{greenbox}
\small
\renewcommand{\arraystretch}{1.35}
\begin{tabularx}{\linewidth}{@{} >{\bfseries}p{3.1cm} X @{}}
Sampling frame & The list of individuals a study can actually choose from. Anyone off the
       frame cannot be sampled. \\
Probability sample & Chance does the choosing; every individual on the frame has a known
       chance of selection. All four techniques today are probability samples. \\
Convenience sample & Whoever is easiest to reach --- the researcher or the schedule, not
       chance, chose. Test: \emph{who could never have been picked?} \\
Simple random sample & Every individual --- and every group of $n$ individuals --- has an
       equal chance. Number the frame, draw, skip out-of-range and repeat values. \\
Systematic sample & From an ordered list, every $k$th individual after a random start, where
       $k = N \div n$. Only the start is random. \\
Stratified sample & The population is split into \emph{strata} whose members are alike
       inside; a random sample is drawn from \emph{every} stratum, in proportion. \\
Cluster sample & The population is split into \emph{clusters}, each meant to be a small mix
       of the whole; a few whole clusters are drawn at random and everyone in them is
       measured. \\
Stratum vs.\ cluster & A stratum is alike \emph{inside} and differs from other strata; a
       cluster is a small mix of the whole. Which technique you may use follows from which
       kind of group you have. \\
\end{tabularx}
\end{skillbox}

% fixedskillbox never splits, so the phase table stays intact on one page.
% THE PHASE TABLE IS A CONTRACT: 5 / 35 / 10 / 10 — the I do is ~20 of the 35 and the Guided
% Practice ~15; the debrief is spoken; the homework start is the second formative read.
\begin{fixedskillbox}[Lesson at a Glance --- \MeetingLength]{frost}
\renewcommand{\arraystretch}{1.25}
\begin{tabularx}{\linewidth}{@{}l c X X@{}}
\textbf{Phase} & \textbf{Min} & \textbf{Students} & \textbf{Teacher} \\
\hline
Warm-Up & 5 & Riverbend's $900$: a percent, that percent of $60$, an interval down a list ---
  silent & Debrief item~3 aloud; write ``$900 \div 60 = 15$; $7$, $22$, $37$'' on the board
  and leave it up --- section~1 names it \\
Guided Notes \& Practice & 35 & Fill the vocabulary box and notes 1--2 as each term is named;
  then \emph{Cedar Ridge Apartments} together, pen in their hand & \textbf{I do} notes 1--2
  (20) $\to$ \textbf{we do} the Guided Practice box (15) \\
Debrief & 10 & Pens down, papers up --- answer aloud, nothing to fill in & Maya's rooms $4$
  and $7$ back on the board; the ``it was random, so it's fine'' reflex, then the
  ``random means unplanned'' reflex; the Bayside cold check \\
Close \& Assign & 10 & \textbf{Start the homework in class}, alone and silent & Announce the
  homework and its form (packet), circulate for the second formative read, preview
  Lesson~1.4 \\
\end{tabularx}
\end{fixedskillbox}

\begin{fixedskillbox}[Warm-Up --- Activate Prior Knowledge (5 min)]{frost}
\begin{tabularx}{\linewidth}{@{}X X@{}}
\begin{minipage}[t]{\linewidth}\vspace{0pt}
    \textbf{The three items and what each seeds}
    \begin{itemize}[leftmargin=1.1em, itemsep=1pt, topsep=2pt]
      \item \textbf{1.} $270$ of $900$ as a percent. \textbf{Spirals} to Lesson 1.1's relative
            frequency; it is the share the strata table in section~2 preserves.
      \item \textbf{2.} $30\%$ of $60$. \textbf{Seeds:} the stratified allocation with the
            name removed --- section~2 opens on this $18$ and asks what they were doing.
      \item \textbf{3.} $900 \div 60 = 15$, start $7$, then $22$ and $37$. \textbf{Seeds:} the
            systematic sample with the name removed --- the second half of section~1 attaches
            the name to arithmetic already done.
    \end{itemize}
\end{minipage}
&
\begin{minipage}[t]{\linewidth}\vspace{0pt}
    \textbf{Running it}
    \begin{itemize}[leftmargin=1.1em, itemsep=1pt, topsep=2pt]
      \item \textbf{Debrief item~3 aloud} and write $900 \div 60 = 15$; $7$, $22$, $37$ on the
            board. Leave it up all period; section~1 hangs \emph{systematic sample} on that line.
      \item Item~3(b) is the tell. Students who write $8$ and $9$ read ``next'' as ``the next
            name'' --- fix that misread at the board in ten seconds before the notes begin.
      \item Say that all three are Algebra~1 arithmetic and are not being retested. \textbf{Do
            not name any technique here} --- the warm-up is deliberately wordless about them.
    \end{itemize}
\end{minipage}
\end{tabularx}
\end{fixedskillbox}

\boxguard
\begin{skillbox}[Hook --- before anyone sits down]{frost}
The hook is on \textbf{page 1 of the notes}, under the vocabulary box, and on the board:
\textit{``The council needs $60$ of Riverbend's $900$ students. Maya puts the $30$ homeroom
numbers in a hat, draws two --- rooms $4$ and $7$ --- and surveys all $30$ students in each.
Pure chance.''} \textbf{Ask: is that as trustworthy as drawing $60$ names out of the hat?}
Students circle \emph{yes}, \emph{no}, or \emph{it depends} and write one line of reason in the
hook box \emph{before} anyone speaks; then take the hand vote and write the split on the
board. Most rooms vote \emph{yes} --- ``it was random.'' \textbf{Do not settle it, and do not
mention that homerooms are grouped by grade.} Say only that the answer is the last thing in
section~2 and that they will vote again then. The vote is what makes the crux land: students
have to have committed to \emph{chance chose, so it is fine} before the notes take it away.
\end{skillbox}

\boxguard
\begin{skillbox}[Guided Notes \& Practice --- I do, then we do (35 min)]{frost}
\textbf{This is the whole lesson.} There is no group activity, and there is no independent
practice set on this page: \textbf{the homework is the individual practice}, started in class.
\begin{multicols}{2}
\textbf{I do (about 20 min).} Two sections, each in two moves; students fill the blanks as each
idea is named, and the vocabulary box on page~1 is filled \emph{as you go}, never in advance.
\textbf{Every fill-in cluster opens with a word bank}; the work is choosing and justifying,
not recalling.

\emph{Section 1} opens on the frame and the two ways to choose, then the four-row \emph{who
chose it?} table. \textbf{Take the third-period row last and take a hand count before anyone
writes}: half the room calls it random because the teacher picked no favorites. Then ask
\emph{could a fourth-period student have been chosen?} and let the red callout close it ---
\emph{random} does not mean \emph{unplanned}. Five minutes. Its second half is the two list
techniques: run the digit strip live and \textbf{stop at $951$} --- vote on what to do with it
before saying anything; the skip is the only non-obvious move in an SRS. Then the systematic
paragraph, which is the warm-up with its name: $k = 15$, start $7$, and \emph{what is actually
random here?} Only the start. Five minutes.

\textbf{Section 2 is the lesson.} Its first half is stratified: the warm-up already computed
$18$ --- start there and ask what they were doing. Fill the $1/15$ table (one division per row),
work the 10th-grade block, and spend the saved minutes on the word \emph{guarantees}. Four
minutes. Its second half, \emph{Cluster Sample --- and Maya's Homeroom Draw}, is the crux.
Define cluster, work $2 \times 30 = 60$, and \textbf{only then reveal the homeroom map}. Fill
the four short blanks with the room --- both 9th-grade rooms, zero seniors, chance chose, nobody
counted wrong --- and \textbf{re-take the hook vote} before saying a word. Then land the blue
callout: a homeroom is alike inside, so it is a stratum; alike inside means \emph{some from
every group}. Close with the contrast table, and make the class say the one-liner aloud before
they write it. Six minutes.

\columnbreak
\textbf{We do (about 15 min).} The Guided Practice box, \emph{Cedar Ridge Apartments}, alone on
page~4 --- the same moves the homework will ask alone, on an apartment complex instead of a
school. Read the building layout aloud when you open it: buildings $1$--$6$ studios, $7$--$20$
one-bedrooms, $21$--$30$ two-bedrooms --- part (c) does not work if students skim past it.
\textbf{Students hold the pen; you hold the questions.} Part (a) goes on them cold, one row at a
time. Part (b) is warm-up arithmetic. \textbf{Part (c) is the crux transferred}: the hat gives
buildings $2$, $5$, and $6$. Take a hand count on \emph{is this sample fine?} before anyone
writes, and make both sides defend it. Part (d) is PS.DC.2b --- the reason, not the word. Part
(e) can be said aloud if time is short.

Circulate during Guided Practice and demand the reason, not the word:
\begin{itemize}[leftmargin=1.1em, nosep]
  \item \textbf{Part (a):} ``How many of the groups did you use --- all of them, or a few?''
        The group names are not evidence.
  \item \textbf{Part (b):} ``You wrote $12$ for the studios. Twelve out of how many?''
  \item \textbf{Part (c):} ``Chance chose the buildings. So what went wrong?'' A student who
        says \emph{bad luck} has not left the hook vote; the answer is that a building is not
        a mix of the complex.
  \item \textbf{Part (d):} ``Finish the sentence --- \emph{because the groups are}\,\ldots''
\end{itemize}
While circulating, pick the papers to put up in the debrief --- one part (c) that named the
building a stratum, and one that said \emph{unlucky draw}. \textbf{Your formative read comes
twice}: here, and again in the last ten minutes as students start the homework.
\end{multicols}
\end{skillbox}

\boxguard[24]
\begin{skillbox}[Debrief --- whole class, spoken (10 min)]{frost}
Pens down, papers up. \textbf{This debrief is spoken} --- there is no box at the end of the
notes to fill in, so it lives entirely on the board and in the room.
\begin{multicols}{2}
\begin{enumerate}[leftmargin=1.3em, nosep, itemsep=3pt]
  \item \textbf{Put Maya's rooms $4$ and $7$ back up} and make a student say the sentence:
        chance chose, nobody counted wrong, and not one senior was asked --- because a homeroom
        is a stratum. Then state it as PS.DC.2b: the technique has to match the groups.
  \item Put up the two Guided Practice part (c)s you picked. The difference between them is
        not arithmetic --- it is whether the student still believes \emph{random} settles it.
  \item Circle the third-period row and take the \emph{random-means-unplanned} reflex:
        \emph{``Who could never have been picked? Then who chose?''}
  \item Take part (d) last, and insist on the form: \emph{stratified, because the groups are
        alike inside and differ on the very thing being asked}.
\end{enumerate}
\columnbreak
\textbf{Check for understanding --- ask it cold, whole class.} \emph{``Bayside Middle School:
$750$ students in three grades, homerooms grouped by grade, $30$ to survey. Plan~1 draws $2$
of the $30$ homerooms and asks everyone in them. Plan~2 takes a random $10$ from each grade.
Name each technique. Which one guarantees all three grades are heard --- and why doesn't the
other, even though it is random?''} Hands down, thirty seconds of think time, then cold-call
three students. Correct answer: Plan~1 is \textbf{cluster}, Plan~2 is \textbf{stratified};
only Plan~2 guarantees it, because chance can land both of Plan~1's rooms in one grade --- a
homeroom is alike inside, so it is a stratum, not a mix of the school. \emph{This replaces the
exit ticket.}

\textbf{Formative read.} Sort the three answers as they come: \textbf{(a)} named both and said
\emph{alike inside}; \textbf{(b)} named both and said Plan~1 ``might get unlucky''; \textbf{(c)}
said Plan~1 is fine because it is random, or swapped the names. Pile (b) is the teachable one
--- it is the right instinct with the wrong lever (the procedure, not luck). \textbf{If pile
(c) runs past a third of the room, open Lesson~1.4 by re-running section~2's second half}, and
say so plainly.

\textbf{If the debrief runs short}, cut items 3 and~4 --- item~1 is the one that cannot be
cut. \textbf{Do not let the debrief eat the last ten minutes} --- the homework start is not
optional padding, it is your second formative read.
\end{multicols}
\end{skillbox}

\boxguard
\begin{skillbox}[Homework --- scored, started in class, due after two study halls]{goldbox}
Six exercises replanning the \textbf{Harbor Point Community Pool} study ($1{,}500$ members,
$100$ to survey; the board threw out 1.2's Tuesday-evening sample): \textbf{1} the
Tuesday-evening study --- population, sample, \emph{convenience}, and who had no chance
(\emph{the spiral to Lesson 1.2}); \textbf{2} Plan A, the generator table, with one
out-of-range value and one repeat --- two different reasons to discard; \textbf{3} Plan B, the
stratified allocation at $1/15$; \textbf{4} Plan C, $k = 15$ from a start of $9$; \textbf{5}
Plans B and D named side by side --- the \textbf{contrast pair} --- then the four sessions
drawn are three lap-swim hours: \emph{the crux head-on}, a random draw of clusters that are
alike inside; \textbf{6} the board's requirement that seniors be heard --- choose the plan and
justify it (PS.DC.2b). Capped at \textbf{two pages} (user decision 2026-09-06); the old
``name the technique'' matching table and the extension are gone.

\textbf{Start it in the last ten minutes of the period, in class and alone.} \textbf{Scoring:}
12 points --- 2 per item. \textbf{Item~5 is where the misconception shows}: a student who
explains the lap-swim draw by ``they got unlucky'' or ``the sample was too small'' has the
wrong lever. \textbf{Item~2 carries the second check}: ``no'' in both discard rows without
distinguishing out-of-range from repeat is half the skill.

\textbf{Packet or Desmos --- decided here.} The packet pages are authored and are the default.
Desmos Classroom has random-sampling simulations, but nothing that asks a student to decide
whether a group is a stratum or a cluster and justify a technique from it --- the items this
lesson exists for --- so \textbf{1.3 is a packet night}. Say so aloud at Close \& Assign.

\textbf{Preview of Lesson~1.4.} Every technique today assumed the people chosen actually
answer, answer honestly, and read the question the same way. Next class drops all three
assumptions: undercoverage, nonresponse, response bias, question wording --- PS.DC.2c and
AFDA.DA.2e--g. Do not open it here.

\textbf{Connections \& Big Ideas.} \textit{Unit 1 core idea:} a conclusion is only as
trustworthy as the study that produced it. \textit{Standards covered:} PS.DC.2a, PS.DC.2b;
AFDA.DA.2b (with PS.DC.1d carried forward). \textit{Spirals forward to:} bias (1.4, PS.DC.2c),
observational studies (1.5, PS.DC.2d), randomization in experimental design (1.6,
PS.DC.3a--b), the full-cycle project (1.8), and --- directly --- sampling distributions and
margin of error (7.1--7.2). \textit{Reaches back to:} Lesson 1.2's population, sample, and
sampling variability, and Algebra~1 percent and proportional reasoning.
\end{skillbox}

\boxguard
\begin{skillbox}[Watch For (while circulating)]{redbox}
\begin{itemize}[leftmargin=1.2em, nosep]
  \item \textbf{``It was random, so it's fine''} (notes~2, GP~(c), HW~5 --- the target
        misconception). Probe: \emph{``Chance chose the rooms. Is a homeroom a mix of the
        school, or one grade? Then what did the draw actually sample?''}
  \item \textbf{Stratified and cluster swapped} (notes~2 table, GP~(a), HW~5). Probe:
        \emph{``How many of the groups did you use --- all of them, or a few?''} All
        $\Rightarrow$ stratified. A few $\Rightarrow$ cluster. The group names are not evidence.
  \item \textbf{``Random'' meaning ``unplanned''} (notes~1 row~3, GP~(a) mailbox row, HW~1).
        Probe: \emph{``Who could never have been picked?''} If that list is non-empty, chance
        did not choose.
  \item \textbf{The interval computed as $n \div N$} (notes~1, HW~4): $60 \div 900$ gives
        $0.0\overline{6}$ and the student moves on. Probe: \emph{``What are the units of $k$
        --- names, or fractions?''}
  \item \textbf{``Skip it'' meaning ``stop''} (notes~1 digit strip, HW~2). The generator keeps
        running; only that number is discarded. HW~2 has two different discard reasons ---
        out of range and repeat --- and both must be named.
  \item \textbf{A technique named with no reason} (GP~(d), HW~6). PS.DC.2b is the
        justification, not the label. Probe: \emph{``Finish the sentence: because the groups
        are\,\ldots''}
\end{itemize}
Cold-call prompts: \emph{``Who chose that sample?''} \quad \emph{``What list do they have?''}
\quad \emph{``Is that group alike inside, or a mix of the whole?''}
\end{skillbox}

\boxguard
\begin{skillbox}[Close \& Assign (10 min)]{goldbox}
\textbf{Students start the homework here, in class, alone and silent --- this is the individual
practice, and the first minutes of it are your best formative read.} Hand the launch first ---
six exercises on paper, scored out of 12, due at the start of the first class after two study
halls --- and point out that the \emph{Keep in Mind} box on the cover is the page to flip back
to. Circulate: \textbf{item~5 is the column that tells you whether section~2 landed}, and it is
on page~2, so put students on 1, 5, and~6 first and let 2, 3, and~4 go home. Sort what you see
into three piles as you walk --- said \emph{a session is alike inside, not a mix} (ready); said
\emph{lap swimmers all answer the same} with no rule (ready, needs the language); said
\emph{unlucky draw} or \emph{too small} (the misconception survived) --- and open Lesson~1.4 by
re-running section~2's second half if that third pile ran past a third of the room. Then close
on what changed: \emph{``You walked in able to divide $900$ by $60$. You walk out knowing that
chance choosing is the price of admission, not the finish line --- the technique has to match
how the groups are built, and a random draw of groups that are alike inside can miss a whole
grade without anyone counting wrong.''} \textbf{Preview:} Lesson~1.4, \emph{Bias in Samples and
Surveys} --- today every chosen student answered; next time we ask what happens when the
frame is short, the chosen refuse, and the question leans.
\end{skillbox}

% ── Teacher notes — one per component, in packet order (three) ──────────────
% Teacher-only prose lives HERE, never in a _key: a note is the one block in a key with no
% counterpart in the blank. No note for the debrief or the homework start — both are fully
% specified in their own boxes above.
\begin{teachernote}[Warm-Up]
Five minutes, silent, timed. All three items are Algebra~1 arithmetic and should be near
universal --- say so when you assign it, because item~3 looks unfamiliar and is not.
\textbf{Item~3 is the one to read while collecting}: it is a systematic sample with the name
removed --- $900 \div 60 = 15$, start at $7$, next are $22$ and $37$. Students who write $22$
and $37$ have already done the second half of section~1; students who write $8$ and $9$ read
``next'' as ``the next name,'' and that misread is fixed at the board in ten seconds before the
notes begin. Item~2's $18$ is the stratified allocation without its name; section~2 opens on it.
Write $900 \div 60 = 15$; $7$, $22$, $37$ on the board and \textbf{leave it up all period}. Do
\emph{not} name any of the four techniques here --- the warm-up is deliberately wordless about
them; the notes name what the class has already computed.
\end{teachernote}

\begin{teachernote}[Guided Notes \& Practice]
\textbf{This one document is the lesson --- pace it 20 / 15, then 10 for the debrief, then 10
for the homework start.} Section~1 must move: five minutes on the frame and the \emph{who
chose it?} table, five on the two list techniques. \textbf{Take the hand count on the
third-period row before anyone writes}; if students fill the row first the trap is gone and the
red callout becomes a definition to copy. Run the digit strip live, not read: stop at $951$ and
vote before saying anything --- the skip is the only non-obvious move in an SRS. The systematic
paragraph is the warm-up with its name; do not re-teach the arithmetic.

\textbf{Spend the time in section~2}, which carries the misconception. The strata table is
one division per row --- four minutes --- and the minute that matters is the word
\emph{guarantees}. Then slow down. Define cluster, work $2 \times 30 = 60$, and \textbf{only
then reveal that homerooms $1$--$9$ are the 9th grade}; if the map is on the board before the
draw, the crux collapses into a definition. Fill the four short blanks with the room ---
9th, $0$, Yes, No --- and \textbf{re-take the hook vote before you say a word}. Let the room
argue; the students who voted \emph{yes} will reach for \emph{bad luck}. Then land the blue
callout: not luck --- \emph{about one draw in four} lands both rooms in one grade, because a
homeroom is alike inside, so it is a stratum. Close the section with the contrast table and the
spoken one-liner: \emph{stratified is some from every group; cluster is everyone from a few}.

\textbf{Guided Practice is where the pen changes hands.} \emph{Cedar Ridge Apartments} runs the
moves the homework will ask alone --- name the plans, allocate the strata, the cluster draw
that samples one stratum, choose and justify --- on an apartment complex instead of a school.
Read the building layout aloud when you open the box. Part~(c) is the crux transferred: a
building holds one unit type, so it is a stratum wearing a cluster's clothes, and buildings
$2$, $5$, $6$ are sixty studio households with the car owners never asked. A student who
defends it with ``but it was random'' has not left the hook vote. If you only get through (a),
(b), (c), and (d), that is the right trade; (e) can be said aloud.

\textbf{There is no independent practice set on this page, and that is deliberate --- the
homework is the individual practice.} Guided Practice is the last thing on the page; the period
ends with students starting the homework in front of you, which is where your second formative
read comes from. Item~5 is the one to watch. If the ``unlucky draw'' pile is more than a third
of the room, \textbf{open Lesson~1.4 by re-running the second half of section~2}, and say so
plainly.

\textbf{Debrief (10 min), spoken.} Maya's rooms back on the board with a student saying
\emph{chance chose, and it still missed, because a homeroom is a stratum}; it is the one move
that cannot be cut. \textbf{Do not let the debrief eat the last ten minutes} --- the homework
start is not optional padding.
\end{teachernote}

\begin{teachernote}[Homework]
\textbf{Scored, 12 points (2 per item), due at the start of the first class after two study
halls.} The ten minutes at the end of the period are a \emph{start}, not a completion: put
students on 1, 5, and~6 while you can still catch a wrong start, and let 2, 3, and~4 go home.
Item~5 is the one to read first when scoring --- the sessions were drawn \emph{at random}, so
the answer cannot be ``they picked badly'' or ``too small''; it has to be that a swim session
is not a small mix of the membership, so four clusters can all lean one way. Ask a student who
reaches for size what would happen if the four sessions held $500$ people. Item~2's table has
two different reasons for discarding a number --- $1587$ is out of range and the second $0412$
is a repeat --- and ``no'' in both rows without distinguishing them is half the skill. Item~6 is
PS.DC.2b: \emph{stratified} with no reason is half credit; the reason is that the senior stratum
is sampled separately, so $25$ seniors are guaranteed a place, which no other plan promises.
Score items 1--4 for correctness and items 5--6 for the \emph{reason}, not the verdict.
\textbf{When you score the page, sort item~5 into three piles} --- named \emph{alike inside}
(ready); said lap swimmers all answer the same with no rule (ready, needs the language); said
unlucky or too small (the misconception survived). If more than a third land in the last pile,
reopen Maya's homerooms in the Lesson~1.4 warm-up debrief rather than letting it ride.

\textbf{Packet is the call for 1.3.} Desmos carries sampling simulations and none of the
stratum-or-cluster reasoning in this set; do not swap it in. The \textbf{spiralbox} hands off to
Lesson~1.4 deliberately: every technique here assumed the chosen members answer, answer
honestly, and read the question the same way. Do not open that here.
\end{teachernote}
\end{document}
